% ============================================================
\subsection{Local-State Persistence and Phase-Hamiltonian Realization}
\label{sec:flagship-quantum}
% ============================================================

The phase lane is the realization of \emph{local-state persistence}: a local
coherent state remains itself through the realized temporal order and carries a
reversible scalar transport law on the smallest faithful cut where that mode
is visible. Phase rigidity, transport rigidity, and scalar closure then
determine the physical phase law class on that realized cut using only the
interfaces inherited from the earlier structural chapters together with the
continuum bridge. The phase generator itself is forced
internally by phase generation and selected-cut energy orthogonality. The
recognizable Schr\"odinger form is obtained by local collapse of the already
forced transport/closure class to its canonical Laplace/potential
representative.

Theorem~\ref{thm:temporal-realization-forbids-primitive-internal-creation}
fixes the admissible source meaning on this carrier. In a closed reversible
phase sector there is no primitive visible creation, so any visible
inhomogeneity must appear as exchange with another visible sector or as
returned residual from the hidden complement on the same coherent bundle.

The relevant compatibility identity on this carrier is phase-compatible
reversible transport: the selected cut is the smallest one on which the forced
phase class, scalar recursive transport datum, and scalar closure class are
faithfully visible. The interval \(J\subseteq\bbR\) introduced below is
therefore not extra primitive data; it is a coordinate model of the realized
causal order on the phase carrier.

\paragraph{Carrier status.}
On the current theorem surface, the phase carrier itself is no longer being
introduced from scratch in this flagship section. By
Theorem~\ref{thm:three-forcing-bridges-from-bedrock-to-primitive-physics}, the
primitive direct lanes are already forced on the selected bundle. The
phase-Hamiltonian package below should therefore be read as the compatibility
and representative surface for the already forced \(1\)-channel, not as an
independent sectoral hypothesis beside the temporal realization law.

\subsubsection{Phase-Hamiltonian Realization}

\begin{definition}[Phase-Hamiltonian realization]
\label{def:phase-hamiltonian-realization}\lean{PhaseLane.HamiltonianRealization}
Let \(J\subseteq \bbR\) be a compact time interval, let \(\mathcal W\) be a
Hilbert space with dense test domain \(\mathcal W_0\subseteq\mathcal W\), and
let
\[
\mathcal F=\{(R_n,h_n,L_n)\}_{n\in\bbN}
\]
be a refinement-compatible realized law family into \(\mathcal W\). A
\emph{phase-Hamiltonian realization} on a visible carrier
\(\mathcal X_{\mathrm{ph}}\subseteq\mathcal W\) consists of:
\begin{enumerate}[label=(\roman*),itemsep=0.2em]
\item for every \(n\), the relevant selected channels of \(R_n\) carry
      selected channel transport systems and belong to admissible realized
      classes satisfying the hypotheses of
      Corollary~\ref{cor:intrinsic-package-implies-governing-datum-rigid};
\item the selected channel families are organized into a
      refinement-compatible family on \(C^1(J,\mathcal W_0)\) with a
      refinement-coherent reconstruction datum whose exact core realization is
      bounded;
\item the induced continuum transport class on \(\mathcal X_{\mathrm{ph}}\) is
      represented by a self-adjoint second-order operator
      \[
      \mathcal T_{\chi}:\mathcal W_0\to\mathcal W;
      \]
\item the induced continuum closure class on \(\mathcal X_{\mathrm{ph}}\) is
      represented by a self-adjoint scalar operator
      \[
      V_{\chi}:\mathcal W_0\to\mathcal W.
      \]
\end{enumerate}
\end{definition}

\subsubsection{Least-Motion Phase Carrier}

\begin{theorem}[Least-motion phase observer]
\label{thm:least-motion-phase-observer}\lean{PhaseLane.leastMotionObserver}
Assume the hypotheses of Definition~\ref{def:phase-hamiltonian-realization}.
Suppose:
\begin{enumerate}[label=(\roman*),itemsep=0.2em]
\item the phase carrier \(\mathcal X_{\mathrm{ph}}\) is faithful,
      closure-admissible, and stable under admissible promotion for the
      realized family;
\item no proper visible subcarrier of \(\mathcal X_{\mathrm{ph}}\) still
      supports a forced phase-carrier class, a scalar recursive transport
      datum, and a scalar closure class at the observed scale;
\item any characteristic observer datum satisfying \emph{(i)}--\emph{(ii)} is
      horizon-preserving equivalent to the datum generated by
      \(\mathcal X_{\mathrm{ph}}\).
\end{enumerate}
Then \(\mathcal X_{\mathrm{ph}}\) is the unique least-motion characteristic
observer for the realized family.
\end{theorem}

(Proof in Appendix~\ref{sec:appendix-carrier}.)

\subsubsection{Induced Reversible Phase Law}

\begin{theorem}[Induced reversible phase law on the selected cut]
\label{thm:induced-phase-hamiltonian-law}\lean{PhaseLane.inducedHamiltonianLaw}
Assume the hypotheses of
Theorem~\ref{thm:least-motion-phase-observer}. Then there exist a
representative \(J_{\chi}\) of the forced phase-carrier class on the selected
cut, a unique induced scalar coefficient \(\kappa_{\chi}\ge 0\), and a
remainder \(\mathsf R_{\mathrm{ph}}\in C(J,\mathcal W)\) such that every
\(\psi\in C^1(J,\mathcal W_0)\) satisfies
\[
J_{\chi}\partial_t\psi
=
\bigl(-\kappa_{\chi}\mathcal T_{\chi}+V_{\chi}\bigr)\psi
+
\mathsf R_{\mathrm{ph}}.
\]
By
Theorem~\ref{thm:temporal-realization-forbids-primitive-internal-creation},
\(\mathsf R_{\mathrm{ph}}\) is an exchange remainder rather than a primitive
source term on the selected cut.
If the selected cut is lossless at leading order, so that
\(\mathsf R_{\mathrm{ph}}=0\), this is the forced reversible phase law on the
selected cut.
\end{theorem}

(Proof in Appendix~\ref{sec:appendix-carrier}.)

The theorem above therefore reduces the phase lane to a carrier-level
representative problem. The selected cut already fixes the reversible phase law
class, the scalar transport coefficient, and the admissible exchange
remainder. The next theorem shows that the phase generator is internally
canonical on the selected cut; the local transport/closure collapse then
identifies the Laplace/potential representative.

\begin{theorem}[Phase generation and exact selected-cut energy split force the canonical complex generator]
\label{thm:phase-generation-and-energy-orthogonality-force-complex-generator}\lean{PhaseLane.EndpointFoundationFrontierData.complexGeneratorSurface}
Assume the hypotheses of
Theorem~\ref{thm:induced-phase-hamiltonian-law}. Assume in addition that, on
each relevant selected phase channel, the realized class on the selected cut
is phase-generated in the sense of
Definition~\ref{def:phase-generated-realized-class}, that the realized family
satisfies a canonical selected Schur bridge on that same selected cut. Then
the forced phase-carrier class on the selected cut admits a unique
pairing-compatible strong representative. In standard complex notation, that
representative is the canonical complex phase generator on
\(\mathcal X_{\mathrm{ph}}\).
\end{theorem}

(Proof in Appendix~\ref{sec:appendix-carrier}.)

\subsubsection{Endpoint Identification}

Writing the selected phase representative and transport datum in standard
wave-mechanics form now separates into two carrier-level steps. The phase
generator is already fixed by
Theorem~\ref{thm:phase-generation-and-energy-orthogonality-force-complex-generator}.
The final identification is the unique local pairing-symmetric
transport/closure representative on the selected cut. When that local endpoint
family assumes the standard Laplace/potential form, the resulting
endpoint law is the classical Schr\"odinger equation
\cite{Schrodinger1926,Stone1932}.

\begin{definition}[Phase local endpoint-rigidity package]
\label{def:phase-local-endpoint-rigidity-package}\lean{PhaseLane.EndpointCollapse}
Assume the hypotheses of
Theorem~\ref{thm:phase-generation-and-energy-orthogonality-force-complex-generator}.
A \emph{phase local endpoint-rigidity package} on the selected phase carrier
\(\mathcal X_{\mathrm{ph}}\) consists of:
\begin{enumerate}[label=(\roman*),itemsep=0.2em]
\item the minimal pairing-symmetric transport family on
      \(\mathcal X_{\mathrm{ph}}\) is represented levelwise by a
      refinement-compatible local scalar stencil family of minimal local order
      \(2\);
\item on each relevant selected phase channel, the corresponding one-step
      transport data are without preferred direction at quadratic order in the
      sense of
      Definition~\ref{def:no-preferred-direction-selected-channel-system};
\item the transport stencil family carries a refinement-coherent
      reconstruction datum on \((\mathcal W,\mathcal W_0)\) whose exact core
      realization is bounded;
\item the scalar closure family on \(\mathcal X_{\mathrm{ph}}\) is represented
      levelwise by a refinement-compatible local scalar stencil family of
      order \(0\);
\item the scalar closure family carries a refinement-coherent reconstruction
      datum on \((\mathcal W,\mathcal W_0)\) whose exact core realization is
      bounded.
\end{enumerate}
\end{definition}

\begin{theorem}[Phase local endpoint collapse to Laplace/potential form]
\label{thm:phase-local-endpoint-collapse}\lean{PhaseLane.endpointCollapse}
Assume the hypotheses of
Definition~\ref{def:phase-local-endpoint-rigidity-package}. Then there exist a
forced positive Laplace-type representative
\[
\Delta_{\chi}:\mathcal W_0\to\mathcal W
\]
for the selected transport class and a forced real multiplication operator
\[
V:\mathcal W_0\to\mathcal W
\]
for the selected scalar closure class such that every
\(\psi\in C^1(J,\mathcal W_0)\) satisfies
\[
i\,\partial_t\psi
=
\bigl(-\kappa_{\chi}\Delta_{\chi}+V\bigr)\psi
+
\mathsf R_{\mathrm{ph}}.
\]
If the selected cut is lossless at leading order, so that
\(\mathsf R_{\mathrm{ph}}=0\), this is exactly the Schr\"odinger equation on
the selected phase carrier.
\end{theorem}

(Proof in Appendix~\ref{sec:appendix-carrier}.)

\begin{corollary}[Recognizable Schr\"odinger identity]
\label{cor:recognizable-schrodinger-identity}\lean{PhaseLane.recognizableIdentity}
Assume the hypotheses of
Definition~\ref{def:phase-local-endpoint-rigidity-package}. Then the selected
phase carrier satisfies the Schr\"odinger identity
\[
i\,\partial_t\psi
=
\bigl(-\kappa_{\chi}\Delta_{\chi}+V\bigr)\psi
+
\mathsf R_{\mathrm{ph}},
\]
where \(\Delta_{\chi}\) is the forced positive Laplace-type transport
representative and \(V\) is the forced real multiplication potential. If the
selected cut is lossless at leading order, so that \(\mathsf R_{\mathrm{ph}}=0\),
this is exactly the Schr\"odinger equation on the selected carrier.
\end{corollary}

\begin{proof}
This is exactly Theorem~\ref{thm:phase-local-endpoint-collapse}.
\end{proof}

\begin{corollary}[Exact Schur form of the selected phase remainder]
\label{cor:phase-schur-remainder-identity}\lean{PhaseLane.schurRemainder}
Assume the hypotheses of
Theorem~\ref{thm:phase-local-endpoint-collapse}. Assume in addition that the
same selected datum satisfies a canonical selected Schur bridge. Then there is
a unique phase-projected Schur correction
\[
\Sigma_{h,\mathrm{ph}}(z)\in C(J,\mathcal W)
\]
obtained by reading the common selected Schur correction \(\CCSigma{h}{z}\)
through the forced phase endpoint identification on the selected carrier, and
\[
\mathsf R_{\mathrm{ph}}=\Sigma_{h,\mathrm{ph}}(z).
\]
Accordingly, every \(\psi\in C^1(J,\mathcal W_0)\) satisfies
\[
i\,\partial_t\psi
=
\bigl(-\kappa_{\chi}\Delta_{\chi}+V\bigr)\psi
+
\Sigma_{h,\mathrm{ph}}(z).
\]
\end{corollary}

(Proof in Appendix~\ref{sec:appendix-carrier}.)

\subsubsection{Concrete Hidden-Channel Phase Model}

The exact Schur form above becomes computationally informative once one
computes the returned phase profile on a concrete realized family. The
simplest solvable case is a visible phase carrier locally coupled to one
hidden phase carrier on the same primitive lattice. On that model the
phase remainder is an exact self-energy term in the usual spectral sense.

\begin{definition}[Coupled visible-shadow phase lattice model]
\label{def:coupled-visible-shadow-phase-lattice-model}\lean{PhaseLane.HiddenChannel.LatticeModel}
Assume the hypotheses of
Corollary~\ref{cor:phase-schur-remainder-identity}. A
\emph{coupled visible-shadow phase lattice model} consists of:
\begin{enumerate}[label=(\roman*),itemsep=0.2em]
\item a one-dimensional primitive lattice with spacing \(\ell_{\chi}>0\);
\item a visible lattice phase law
      \[
      i\,\partial_t\psi
      =
      \bigl(-\kappa_{\chi}\Delta_{\chi}+U_{\chi}\bigr)\psi
      +
      \gamma_{\chi}\eta,
      \]
      where
      \[
      (\Delta_{\chi}\psi)_n
      =
      \frac{\psi_{n+1}-2\psi_n+\psi_{n-1}}{\ell_{\chi}^2},
      \qquad
      U_{\chi}\in\bbR;
      \]
\item a hidden lattice phase law
      \[
      i\,\partial_t\eta
      =
      \bigl(-\widetilde\kappa_{\chi}\Delta_{\chi}+\widetilde U_{\chi}\bigr)\eta
      +
      \gamma_{\chi}\psi,
      \]
      with \(\widetilde\kappa_{\chi}\ge 0\),
      \(\widetilde U_{\chi}\in\bbR\), and \(\gamma_{\chi}\ge 0\);
\item a common selected plane-wave sector on which
      \[
      \psi_n(t)=\Psi\,e^{i(nk\ell_{\chi}-\omega t)},
      \qquad
      \eta_n(t)=H\,e^{i(nk\ell_{\chi}-\omega t)}
      \]
      for amplitudes \(\Psi,H\in\mathbb C\).
\end{enumerate}
\end{definition}

\begin{theorem}[Exact hidden-channel phase remainder and spectral splitting]
\label{thm:hidden-channel-phase-remainder-and-spectral-splitting}\lean{PhaseLane.HiddenChannel.splitting}
Assume the hypotheses of
Definition~\ref{def:coupled-visible-shadow-phase-lattice-model}. Define the
visible and hidden phase profiles
\[
E_{\mathrm{vis}}(k)
:=
U_{\chi}
+
\frac{4\kappa_{\chi}}{\ell_{\chi}^2}
\sin^2\!\left(\frac{k\ell_{\chi}}{2}\right),
\]
\[
E_{\mathrm{hid}}(k)
:=
\widetilde U_{\chi}
+
\frac{4\widetilde\kappa_{\chi}}{\ell_{\chi}^2}
\sin^2\!\left(\frac{k\ell_{\chi}}{2}\right).
\]
Then on every mode with \(E_{\mathrm{hid}}(k)\neq \omega\), exact elimination
of the hidden amplitude gives the returned phase profile
\[
\sigma_{\chi}^{\mathrm{ph,hid}}(k,\omega)
=
\frac{\gamma_{\chi}^2}{E_{\mathrm{hid}}(k)-\omega}.
\]
Accordingly the visible phase law is
\[
\omega
=
E_{\mathrm{vis}}(k)
+
\sigma_{\chi}^{\mathrm{ph,hid}}(k,\omega),
\]
equivalently
\[
(\omega-E_{\mathrm{vis}}(k))
(\omega-E_{\mathrm{hid}}(k))
=
\gamma_{\chi}^2.
\]
Hence the two exact mixed phase branches are
\[
\omega_{\pm}(k)
=
\frac{
E_{\mathrm{vis}}(k)+E_{\mathrm{hid}}(k)
\pm
\sqrt{
\bigl(E_{\mathrm{vis}}(k)-E_{\mathrm{hid}}(k)\bigr)^2
+
4\gamma_{\chi}^2
}
}{2}.
\]
\end{theorem}

(Proof in Appendix~\ref{sec:appendix-carrier}.)

\begin{corollary}[Exact off-resonant phase line shifts]
\label{cor:exact-off-resonant-phase-line-shifts}\lean{PhaseLane.HiddenChannel.offResonantShifts}
Assume the hypotheses of
Theorem~\ref{thm:hidden-channel-phase-remainder-and-spectral-splitting}. Define
\[
\delta_{\chi}^{\mathrm{ph}}(k):=E_{\mathrm{hid}}(k)-E_{\mathrm{vis}}(k).
\]
If \(\delta_{\chi}^{\mathrm{ph}}(k)>0\), then
\[
\omega_-(k)
=
E_{\mathrm{vis}}(k)
-
\frac{
2\gamma_{\chi}^2
}{
\delta_{\chi}^{\mathrm{ph}}(k)
+
\sqrt{
\bigl(\delta_{\chi}^{\mathrm{ph}}(k)\bigr)^2+4\gamma_{\chi}^2
}
},
\]
\[
\omega_+(k)
=
E_{\mathrm{hid}}(k)
+
\frac{
2\gamma_{\chi}^2
}{
\delta_{\chi}^{\mathrm{ph}}(k)
+
\sqrt{
\bigl(\delta_{\chi}^{\mathrm{ph}}(k)\bigr)^2+4\gamma_{\chi}^2
}
}.
\]
In particular, away from hidden resonance the visible-like phase branch is
shifted downward and the hidden-like branch upward.
\end{corollary}

(Proof in Appendix~\ref{sec:appendix-carrier}.)

\begin{corollary}[Exact avoided crossing of the coupled phase branches]
\label{cor:exact-avoided-crossing-coupled-phase-branches}\lean{PhaseLane.HiddenChannel.avoidedCrossing}
Assume the hypotheses of
Theorem~\ref{thm:hidden-channel-phase-remainder-and-spectral-splitting}. Suppose
there exists a selected mode \(k_\ast\) such that
\[
E_{\mathrm{vis}}(k_\ast)=E_{\mathrm{hid}}(k_\ast)=:E_\ast.
\]
Then
\[
\omega_{\pm}(k_\ast)=E_\ast\pm\gamma_{\chi},
\]
so the exact spectral gap at the would-be crossing is
\[
\omega_+(k_\ast)-\omega_-(k_\ast)=2\gamma_{\chi}.
\]
In particular, if \(\gamma_{\chi}>0\), the two phase branches do not cross.
\end{corollary}

(Proof in Appendix~\ref{sec:appendix-carrier}.)

\begin{corollary}[Low-energy renormalized phase offset and transport coefficient]
\label{cor:low-energy-renormalized-phase-offset-and-transport-coefficient}\lean{PhaseLane.HiddenChannel.lowEnergyRenormalization}
Assume the hypotheses of
Theorem~\ref{thm:hidden-channel-phase-remainder-and-spectral-splitting}. Define
\[
\Delta_0^{\mathrm{ph}}:=\widetilde U_{\chi}-U_{\chi},
\qquad
\Delta_1^{\mathrm{ph}}:=\widetilde\kappa_{\chi}-\kappa_{\chi},
\]
and assume \(\Delta_0^{\mathrm{ph}}>0\). Then as \(k\to 0\), the lower exact
phase branch has the expansion
\[
\omega_-(k)
=
U_{\chi,\mathrm{eff}}
+
\kappa_{\chi,\mathrm{eff}}\,k^2
+
O(k^4),
\]
where
\[
U_{\chi,\mathrm{eff}}
=
\frac{
U_{\chi}+\widetilde U_{\chi}
-
\sqrt{
\bigl(\Delta_0^{\mathrm{ph}}\bigr)^2+4\gamma_{\chi}^2
}
}{2},
\]
\[
\kappa_{\chi,\mathrm{eff}}
=
\frac{
\kappa_{\chi}+\widetilde\kappa_{\chi}
-
\frac{
\Delta_0^{\mathrm{ph}}\Delta_1^{\mathrm{ph}}
}{
\sqrt{
\bigl(\Delta_0^{\mathrm{ph}}\bigr)^2+4\gamma_{\chi}^2
}
}
}{2}.
\]
Thus hidden-channel return renormalizes both the phase offset and the long-wave
transport coefficient on the visible-like branch.
\end{corollary}

(Proof in Appendix~\ref{sec:appendix-carrier}.)

The two-channel phase model gives an exact coherent realization of the returned
phase profile. On this model, the remainder manifests as spectral line shifts,
avoided crossing, and low-energy renormalization of the visible phase offset
and transport coefficient. Irreversible decoherence would require a richer
hidden-sector family, such as a continuum or many-channel return model, and is
not part of the present two-channel calculation.
